\documentclass{article}
\usepackage[utf8x]{inputenc}
\usepackage[english,hebrew]{babel}
%\selectlanguage{English}
\usepackage{float} 
\usepackage[top=2cm,bottom=2cm,left=2.5cm,right=2cm]{geometry}
\usepackage{graphicx}
\graphicspath{ {./images/} }
\usepackage{fancybox} %fancy oval thing
\usepackage{lipsum} %fancy oval thing
\usepackage{amsmath}



%Fancy oval thing on the first page only

% full page
\thisfancyput(3.25in,-4.5in){%
  \setlength{\unitlength}{1in}\fancyoval(8,10)}%


% smaller oval thing
%\thisfancyput(3.25in,-5.5in){%
%  \setlength{\unitlength}{1in}\fancyoval(8,7)}%



\begin{document}



\selectlanguage{English}
\begin{titlepage}
    \begin{center}
        \vspace*{1cm}
        
        \begin{figure}[t]
        \centering
        \includegraphics[width=7cm]{mb.png}
        \end{figure}

        \vspace{2cm}
        
        \Large
        
        \selectlanguage{English} \textbf{Research Proposal for a}  \hfill
        \selectlanguage{Hebrew} \text{תכנית מחקר לעבודת גמר}
        
        \selectlanguage{English} \textbf{Master of Science Thesis}  \hfill
        \selectlanguage{Hebrew} \text{לקראת תואר מוסמך למדעים}
        
        
        \Large
        \vspace{1.5cm}
     
        \selectlanguage{English} \textbf{By}  \hfill
        \selectlanguage{Hebrew} \text{מאת}
        
        \selectlanguage{English} \textbf{Name here}  \hfill
        \selectlanguage{Hebrew} \text{שם פה}
        
        %\vspace{0.5cm}
        \LARGE
        
     
        \vspace{3cm}
      
        \selectlanguage{English} \textbf{Title}  
        
        \selectlanguage{Hebrew} \text{הנושא}
        
        
        \vspace{3.5cm}
        
        \large
        \selectlanguage{English} \textbf{PI name}  \hfill
        \selectlanguage{Hebrew} \text{מנחה}
        
       
        \vspace{0.5cm}
            
       
        \Large
        \selectlanguage{English} \textbf{October 2020} \hfill
        \selectlanguage{Hebrew} \text{אוקטובר תשפ"א}
        
            
    \end{center}
\end{titlepage}



\selectlanguage{English}




\begin{abstract}

\large
    
Write here your abstract! Up to half a page in length. Make sure to cover the main ideas of you proposal


\end{abstract}

\large
\vspace{1.2cm}

\section{Introduction and Preliminary Results}

Write here your introduction and preliminary results.
This includes scientific background and definition of the open issues in your field. You may include preliminary results by others in your group, as well as results you have already obtained that are relevant to your project.
Now I will quote the first paper \cite{paper1} like this.


\begin{figure}[H]
\centering
\includegraphics[width=8cm]{cat.jpeg}
\caption{Here goes a caption for this picture: this is a cat}
\label{fig:f1}
\end{figure}



\section{Goals}

This section should explicitly state the scientific question you wish to answer, or the hypothesis you wish to examine. It may be presented as an itemized list. Here I will refer Figure~\ref{fig:f1} like this.



\section{Work Plan}

This section should outline the experiments or lines of research you wish to perform one by one, as well as the expected results according to your hypothesis. You should not go into great detail and should not explain standard methodologies. Instead, outline the technical approach, the controls, and possible pitfalls. It should be written in such a way that it will be clear to another scientist outside of your research group.Here I will refer Figure~\ref{fig:f2} like this and also will cite the second paper \cite{paper1}, like so. 



\begin{figure}[H]
\centering
\includegraphics[width=10cm]{random_cats.png}
\caption{Another caption: more cats}
\label{fig:f2}
\end{figure}





\medskip




\begin{thebibliography}{9}

\bibitem{paper1} 
Name of person who wrote paper. 
\textit{"Name of paper"}, 2056.

\bibitem{paper2} 
Name of person who wrote other paper. 
\textit{"Name of other paper"}, 2021.




\end{thebibliography}







\end{document}